\section{投后不是“长期持有”，而是持续再承保}

私募机构拥有投后团队、董事会信息和经营会；个人不具备这些渠道，所以投后流程必须更少、更公开、更可重复。最有效的做法不是全天盯盘，而是把每一项持有理由映射到一个公告、财报或项目运营指标，并在固定时点重新比较事实与原假设。

\begin{exhibitbox}[Exhibit 6：公开信息下的监测节奏]
\small
\begin{tabularx}{\exhibitboxwidth}{L{2.3cm} L{3.5cm} X L{3.4cm}}
\toprule
\textbf{频率} & \textbf{看什么} & \textbf{问题} & \textbf{可能行动} \\
\midrule
月度 & 重大公告、仓位、成交和经济因子重叠。 & 是否出现控制权、审计、融资、监管、停牌或流动性变化？ & 暂停新增；把异常提交红队复核。 \\
季报 / 半年报 & 核心关键经营指标（KPI）、现金流、应收/库存、利润率、债务与再融资。 & 价值桥是否按原口径兑现？熊市假设需要下调吗？ & 维持、减少或退出；不因价格更低而自动加仓。 \\
半年 & 估值、可替代机会和反方论证。 & 基准情景仍有足够余量吗？是否把已经发生的改善重复计价？ & 重新承保；去除无证据的上行。 \\
年度 & 能力圈与方法本身。 & 研究流程是否真的比简单分散化改善决策质量？ & 缩小或归零研究卫星，若过程不能持续执行。 \\
\bottomrule
\end{tabularx}
{\setlength{\emergencystretch}{2em}\sourcenote{本报告的流程框架；个券关键经营指标（KPI）必须来自其自身公开披露。}}
\end{exhibitbox}

\section{硬失效与软失效}

\begingroup
\setlength{\tabcolsep}{4pt}
\begin{tabularx}{\linewidth}{L{2.3cm}L{5.8cm}X}
\toprule
\textbf{失效类型} & \textbf{典型事实} & \textbf{正确反应} \\
\midrule
硬失效 & 财务造假、无法解释的非标审计、控制权或治理崩塌、核心业务/牌照永久受损、债务或再融资击穿熊市假设。 & 立即停止新增；核验事实并在实际流动性、交易规则和价格限制下规划退出。不要假设跌停日一定可以成交。 \\
软失效 & 连续报告期低于原熊市 KPI、预设催化期满仍无兑现、估值余量消失、与既有持仓出现未预料的高重叠。 & 重做价值桥与反方；若无法恢复可验证的上行，减至目标或退出。 \\
价格波动本身 & 一日或一段时间的上涨/下跌，但没有新增基本面事实。 & 触发下一交易日或下一个工作日的研究复核，而非自动把价格变化解释成证据。 \\
\bottomrule
\end{tabularx}
\endgroup

\section{公共关键经营指标（KPI）的选择}

KPI 必须是能把经营叙事和财务结果连接起来的桥，而不是越多越好。制造业可用交付、价格、产能利用率或产品组合；消费/服务可用同店、订单、回款或客户留存；资产运营可用入住率、租约、租户集中度、车流或利用率。每个对象应只保留少量最关键、且来自正式披露的指标。若公司没有披露足以验证核心假设的数据，结论不是去猜数据，而是降低研究等级或放弃。

\section{退出纪律与可选性}

在私募中，退出是预先规划、可能持续多年的过程。个人账户的对应物是：在建仓之前就知道什么事实会让自己不再愿意持有，并保留足够的流动性和未投入资本来处理不确定性。现金在本框架中不是“闲置”，而是避免被迫卖出、避免用杠杆修补错误、等待更好公开证据的可选性。它不保证收益，却是个人最能与机构长期资本相对应的风险管理工具。
